\begin{table}[H]
    \centering
    \caption{\textit{Iris} pradinių duomenų keli pavyzdžiai (12 eilučių – po 4 iš kiekvienos klasės).}
    \begin{tabular}{|c|c|c|c|c|}
        \hline
        \textit{sepal length} & \textit{sepal width} & \textit{petal length} & \textit{petal width} & \textit{class} \\
        \hline
        $5{,}4$ & $3{,}7$ & $1{,}5$ & $0{,}2$ & Iris-setosa \\ \hline
        $4{,}8$ & $3{,}4$ & $1{,}6$ & $0{,}2$ & Iris-setosa \\ \hline
        $4{,}8$ & $3{,}0$ & $1{,}4$ & $0{,}1$ & Iris-setosa \\ \hline
        $4{,}3$ & $3{,}0$ & $1{,}1$ & $0{,}1$ & Iris-setosa \\ \hline
        $7{,}0$ & $3{,}2$ & $4{,}7$ & $1{,}4$ & Iris-versicolor \\ \hline
        $6{,}4$ & $3{,}2$ & $4{,}5$ & $1{,}5$ & Iris-versicolor \\ \hline
        $6{,}9$ & $3{,}1$ & $4{,}9$ & $1{,}5$ & Iris-versicolor \\ \hline
        $5{,}5$ & $2{,}3$ & $4{,}0$ & $1{,}3$ & Iris-versicolor \\ \hline
        $6{,}3$ & $2{,}5$ & $5{,}0$ & $1{,}9$ & Iris-virginica \\ \hline
        $6{,}5$ & $3{,}0$ & $5{,}2$ & $2{,}0$ & Iris-virginica \\ \hline
        $6{,}2$ & $3{,}4$ & $5{,}4$ & $2{,}3$ & Iris-virginica \\ \hline
        $5{,}9$ & $3{,}0$ & $5{,}1$ & $1{,}8$ & Iris-virginica \\ \hline
    \end{tabular}
    \label{tab:iris_pradiniai_duomenys}
\end{table}